\section{Introduction}
% What we do
In this paper, we introduce Grounding-DINO-1.5, a powerful and practical open-set object detection model developed by IDEA-Research. 
% motivation
Object detection serves as a fundamental task in visual recognition, with recent efforts concentrating on creating detectors capable of performing reliably across a variety of real-world environments. 
A key strategy for improving model generalization across diverse object categories is the integration of language data, which has garnered significant attention and has undergone extensive development in recent research.


% background
Grounding-DINO~\cite{GroundingDINO} represents a significant advancement in this area. 
Building on the Transformer-based DINO~\cite{DINO} architecture, it incorporates linguistic information to deliver exceptional performance in various scenarios. 
Following GLIP~\cite{GLIP}, Grounding-DINO redefines object detection as a phrase grounding task, facilitating large-scale pre-training using both detection and grounding datasets. 
This approach, coupled with self-training on pseudo-labeled grounding data derived from an almost limitless pool of image-text pairs, enhances the model’s applicability to open-world settings due to its robust architecture and expansive dataset.

% What we do
Building on the achievements of Grounding-DINO, Grounding-DINO-1.5 expands both the model and data scale, aiming to forge a more powerful and applicable open-set object detection model. 
Specifically, we have increased the model’s size by integrating the pre-trained ViT-L~\cite{EVA-02} model, a more potent vision backbone, and have developed a data engine capable of generating over 20 million large-scale grounding data from different data sources to enrich semantic understanding. 
These enhancements ensure that Grounding-DINO-1.5 aligns more closely with human preferences.

% Eval
Our evaluation extends beyond traditional public benchmarks to include comprehensive human assessments. 
Although many models perform well on standard benchmarks, they often produce less satisfactory visual outputs, primarily due to the disparity between the common metric mean average precision (mAP) and practical real-world application, which typically employs a threshold to filter model outputs. 
The challenge of annotating ground truth labels in open-set scenarios, where consensus is often lacking, further complicates evaluations.
To address these issues, we have instituted a human preference study that compares model outputs directly.

% results
The extensive results from our tests confirm the superiority of Grounding-DINO-1.5, underscoring its enhanced capabilities and practicality in real-world applications.
Specifically, Grounding-DINO-1.5 model achieved 54.3 AP on the COCO detection zero-shot transfer benchmark and simultaneously achieved 55.7 AP and 47.6 on the LVIS-minival and LVIS-val zero-shot transfer benchmarks respectively.



  
  \vspace{-2mm}